\documentclass[11pt]{article}
\usepackage[utf8]{inputenc}
\usepackage{amsmath, amssymb, amsthm}
\usepackage[margin=1in]{geometry}

\theoremstyle{definition}
\newtheorem{definition}{Definition}
\newtheorem{theorem}{Theorem}
\newtheorem{lemma}{Lemma}

\theoremstyle{remark}
\newtheorem{remark}{Remark}

\DeclareMathOperator{\tr}{tr}
\newcommand{\dual}[2]{{}_{V'}\langle #1, #2 \rangle_{V}}

\title{Lecture 11: Compact Embeddings and the Non-Monotone Case}
\date{25.12.2025}
\author{Tien Anh Vu}

\begin{document}
\maketitle

Welcome to Lecture 11 on Stochastic Partial Differential Equations. Please settle down. Today marks a profound turning point in our theoretical framework.

In our previous lectures, we successfully established the existence and uniqueness of SPDEs using the variational approach, relying heavily on the Minty--Browder trick. However, that trick strictly required our operators to be monotone. Today, we ask a vital question: what if our system is not monotone? How do we pass to the limit in our non-linear terms?

The answer, as outlined on the first page of today's notes, is to introduce compact embedding. We now walk through this page step-by-step.

\section*{1. The Main Additional Assumption: Compact Embedding}
We continue to study our core SPDE on the central Hilbert space $H$:
\[
  du_t = A(u_t)\,dt + B(u_t)\,dW_t.
\]

We are still operating within the framework of our Gelfand triple,
\[
  V \hookrightarrow H \simeq H' \hookrightarrow V'.
\]
However, there is now a main additional assumption: the highly regular Banach space $V$ is compactly embedded into $H$. This is denoted by
\[
  V \hookrightarrow \hookrightarrow H.
\]

Analytically, this means that any set that is bounded in the strong topology of $V$ is relatively compact in the strong topology of $H$. Stated formally, if we take a closed ball of radius $K$ in the space $V$, denoted by
\[
  B_K^V(0)=\{v\in V:\|v\|_V\le K\},
\]
then this entire set is a compact subset when viewed inside $H$.

Consequently, if we have a sequence of functions uniformly bounded in $V$, then we are guaranteed to find a subsequence that converges strongly, and not merely weakly, in $H$.

\section*{2. Examples and Geometric Intuition}
To ground this abstraction, let us look at the standard Sobolev spaces. A classical example of a compact embedding is
\[
  H^{1,2}(\Omega) \hookrightarrow \hookrightarrow L^2(\Omega),
\]
provided that the spatial domain $\Omega\subset \mathbb{R}^d$ is bounded.

Why is the boundedness of the domain $\Omega$ required? Suppose instead that our domain is the entire unbounded space $\mathbb{R}^d$. Consider a simple smooth bump function $e(x)$. Now form the sequence of translates
\[
  e(x-n), \qquad n\ge 1.
\]
Because we are only sliding the function along the axis, the spatial norms remain invariant:
\[
  \|e(x-n)\|_{H^{1,2}} = \|e(x)\|_{H^{1,2}}.
\]
Thus this sequence of translated bumps is uniformly bounded in $H^{1,2}$. In particular, each function still has the same finite norm, and the sequence is uniformly bounded in $L^2$ as well. However, as $n\to\infty$, the bump drifts off to positive infinity. It escapes every fixed compact region, so it cannot possibly possess a strongly convergent subsequence in $L^2$.

This is exactly the obstruction to compactness: if one can find a bounded sequence with no strongly convergent subsequence, then the embedding is not compact. Compactness therefore fails on unbounded domains. Translations to infinity destroy compact embeddings.
\[
  H^{1,2}(\mathbb{R}^d) \hookrightarrow \hookrightarrow L^2(\mathbb{R}^d)
\]
is therefore not compact.

\section*{3. Spectral Illustration: The Dirichlet Laplacian}
To give a clearer picture of how compact embeddings work, consider Fourier series on the unit interval and the space $H_0^{1,2}(0,1)$ with Dirichlet boundary conditions.

Consider the Dirichlet Laplacian operator $-\Delta_D$. Its eigenvectors form a complete orthonormal system for the space, given by
\[
  e_k(x)=\sin(k\pi x), \qquad k=1,2,\dots.
\]
Any function $u$ in this space can be represented by its Fourier coefficients
\[
  u_k=\int_0^1 u(x)e_k(x)\,dx.
\]

Crucially, the $H^1$-norm in this spectral representation penalizes high frequencies. This means that oscillatory modes cost more in the $H^1$-norm. In the spectral picture,
\[
  \|u\|_{H^1}^2 \sim \sum_{k=1}^{\infty} (1+k^2)|u_k|^2.
\]
So higher frequencies $k$ are multiplied by a larger weight $k^2$. Low modes are cheap, whereas high modes are expensive. If $u$ contains a lot of high-frequency oscillation, then its $H^1$-norm becomes large. This is exactly what it means to say that the norm penalizes high frequencies: the norm suppresses oscillatory behavior by assigning a large cost to high-frequency coefficients.

A function $u$ belongs to $H_0^{1,2}(0,1)$ if and only if
\[
  \sum_{k=1}^{\infty} k^2 u_k^2 < \infty.
\]

Now consider a bounded set in $V$, meaning we restrict ourselves to functions such that
\[
  \sum_{k=1}^{\infty} k^2 u_k^2 \le M
\]
for some constant $M$. Because of the factor $k^2$, the higher modes are forced to become small as $k\to\infty$. Indeed, for each fixed $k$ one has
\[
  k^2 u_k^2 \le M,
\]
hence
\[
  |u_k| \le \frac{\sqrt{M}}{k}.
\]
So the coefficients decay at least like $1/k$. More importantly, the entire high-frequency tail is small in the $H$-norm. Indeed, for every $N\ge 1$,
\[
  \sum_{k>N} u_k^2
  \le \frac{1}{N^2}\sum_{k>N} k^2 u_k^2
  \le \frac{M}{N^2}.
\]
Hence
\[
  \sum_{k>N} u_k^2 \to 0
  \qquad \text{uniformly as } N\to\infty.
\]
This is the relevant compactness statement: the high-frequency tail becomes negligible in $H$, so most of the $H$-mass is trapped in the first few modes.

Because the higher modes are suppressed, the problem is effectively reduced to finite dimensions. From there, classical real analysis takes over: the Bolzano--Weierstrass theorem guarantees that any bounded sequence in a finite-dimensional space has a convergent subsequence. In this way, one obtains a strongly convergent subsequence in $H$.

\section*{4. Adjusting our Core Assumptions}
At the bottom of the page, we detail how compact embedding affects the core assumptions established earlier in the lecture notes.

What do we keep?

Assumption (V1) --- Coercivity: we retain the coercivity condition
\[
  2\langle A(u),u\rangle + \|B(u)\circ \sqrt{Q}\|_{L_2(U,H)}^2
  \le -\alpha \|u\|_V^2 + \lambda \|u\|_H^2.
\]
This remains the engine of the a priori estimates. It guarantees the uniform boundedness
\[
  \sup_n \mathbb{E}\Bigl(\sup_{t\in[0,T]}\|u_n(t)\|_H^2 + \int_0^T \|u_n(s)\|_V^2\,ds\Bigr)<\infty.
\]
This bound provides the tightness of the Galerkin approximations and ensures that the associated probability mass does not escape to infinity.

\subsection*{4.1. Tightness of the Galerkin Approximations}
If $\mu_n=\mathrm{Law}(u_n)$ denotes the law of the Galerkin approximation $u_n$, then tightness means that these probability measures do not drift off to infinity in the ambient function space. Concretely, for every $\varepsilon>0$, there exists a compact set $K$ such that
\[
  \mu_n(K^c)<\varepsilon
  \qquad \text{for all } n.
\]
Here $K^c$ denotes the complement of $K$ in the ambient function space $X$ where the laws $\mu_n$ are defined, that is,
\[
  K^c=X\setminus K.
\]
Thus, uniformly in $n$, almost all of the probability mass is contained in one fixed compact region. This is the probabilistic compactness statement: after passing to a subsequence, one obtains a convergent subsequence of probability laws, which then provides a candidate limit solution.

We continue by keeping Assumption (V3), namely the linear growth condition on the drift operator.

Assumption (V3) --- Linear Growth:
\[
  \|A(u)\|_{V'}\le M(1+\|u\|_V).
\]
This means that the drift $A(u)$ grows at most linearly with respect to the size of the solution $u$.

Assumption (V5) --- Diffusion Growth: we assume the slightly sub-linear growth condition
\[
  \|B(u)\circ \sqrt{Q}\|_{L_2(U,H)} \le M\bigl(1+\|u\|_V^{1-\delta}\bigr)
\]
for some constants $M,\delta>0$.

What do we drop?

Here lies the central point of today's setup: we completely drop the monotonicity assumption.

By relying on the compact embedding
\[
  V \hookrightarrow \hookrightarrow H,
\]
we will extract strongly convergent subsequences of our Galerkin approximations. Strong convergence allows us to pass to the limit directly inside highly non-linear, non-monotone drift operators $A(u_n)$. We no longer need the algebraic structure of the Minty--Browder trick because topology now performs the essential work.

\section*{5. Finalizing the Assumptions}
We now record the remaining assumptions required to replace the Minty--Browder monotonicity trick with compact embedding.

Assumption (V.6) records the compact embedding already discussed in detail:
\[
  V \hookrightarrow \hookrightarrow H.
\]

Here $V_{\mathrm{weak}}$ means the space $V$ equipped with its weak topology. Thus
\[
  u_n \to u \text{ in } V_{\mathrm{weak}}
\]
means exactly that
\[
  u_n \rightharpoonup u \text{ weakly in } V.
\]
Similarly, $V'_{\mathrm{weak}}$ denotes the dual space $V'$ equipped with its weak topology.

Assumption (V.7) specifies the continuity required once monotonicity has been dropped. We require:
\begin{itemize}
\item the drift operator to be a continuous map
\[
  A:V_{\mathrm{weak}}\cap H \to V'_{\mathrm{weak}},
\]
\item the diffusion operator to be a continuous map
\[
  B:V_{\mathrm{weak}}\cap H \to L_2(U,H).
\]
\end{itemize}

More explicitly, this means the following: if
\[
  u_n \rightharpoonup u \text{ weakly in } V
  \qquad \text{and} \qquad
  u_n \to u \text{ strongly in } H,
\]
then
\[
  A(u_n)\rightharpoonup A(u) \text{ weakly in } V'
\]
and
\[
  B(u_n)\to B(u) \text{ in } L_2(U,H).
\]
Thus the drift operator is continuous in this mixed weak/strong sense, while the diffusion operator is continuous into the Hilbert--Schmidt space equipped with its usual norm topology.

\section*{6. The Topological Bridge: Skorokhod Representation Theorem}
Because we are now relying on compactness, specifically through Prokhorov's theorem, the Galerkin approximations will first yield only weak convergence of probability laws, that is, convergence in distribution.

\subsection*{6.1. Detour: Prokhorov's Theorem and Weak Convergence of Probability Measures}
Prokhorov's theorem is the result that connects tightness with compactness of probability measures. In its standard general form, it is stated on a Polish space, that is, a complete separable metric space.

\begin{theorem}[Prokhorov's theorem]
Let $S$ be a Polish space, and let $\mathcal M\subset \mathcal P(S)$ be a family of probability measures on $S$.
If $\mathcal M$ is tight, meaning that for every $\varepsilon>0$ there exists a compact set $K_\varepsilon\subset S$ such that
\[
  \mu(K_\varepsilon)\ge 1-\varepsilon
  \qquad \text{for all } \mu\in\mathcal M,
\]
then the family of probability measures $\mathcal M$ is relatively compact for the topology of weak convergence on $\mathcal P(S)$.
\par
In particular, if $(\mu_n)$ is a tight sequence of probability measures on $S$, then $(\mu_n)$ is relatively compact in the topology of weak convergence on $\mathcal P(S)$. Hence there exists a subsequence $(\mu_{n_k})$ and a probability measure $\mu$ on $S$ such that
\[
  \mu_{n_k}\Rightarrow \mu.
\]
\end{theorem}

Here the symbol
\[
  \Rightarrow
\]
means weak convergence of probability measures. This is also called convergence in distribution or weak convergence of laws.

Concretely, the convergence
\[
  \mu_{n_k}\Rightarrow \mu
\]
means that for every bounded continuous test function $F$,
\[
  \int F\,d\mu_{n_k}\to \int F\,d\mu.
\]

In the present setting, if
\[
  \mu_n=\mathrm{Law}(u_n),
\]
then Prokhorov's theorem gives a subsequence of the laws of the Galerkin approximations that converges weakly as probability measures.

However, convergence in the weak topology on probability measures is not enough to pass to the limit inside highly non-linear operators such as $A(\cdot)$ and $B(\cdot)$. To do that, one needs almost sure convergence of the random paths themselves.

\subsection*{6.2. Detour: Statement of the Skorokhod Representation Theorem}
The theorem that bridges this gap is the \textbf{Skorokhod representation theorem}. In its standard general form, it is stated on a Polish space, that is, a complete separable metric space.

\begin{theorem}[Skorokhod representation theorem]
Let $S$ be a Polish space, and let $(\mu_n)_{n\ge 1}$ and $\mu$ be probability measures on $S$ such that
\[
  \mu_n \Rightarrow \mu.
\]
Then there exist a probability space $(\widetilde\Omega,\widetilde{\mathcal F},\widetilde P)$ and $S$-valued random variables
\[
  \widetilde X_n,\widetilde X:\widetilde\Omega\to S
\]
such that
\[
  \mathrm{Law}(\widetilde X_n)=\mu_n,
  \qquad
  \mathrm{Law}(\widetilde X)=\mu,
\]
and
\[
  \widetilde X_n \to \widetilde X
  \qquad \text{almost surely on } \widetilde\Omega.
\]
\end{theorem}

Equivalently, if $X_n$ converges in distribution to $X$ as $S$-valued random variables, then one can find copies of these random variables on a new probability space that converge almost surely while preserving exactly the same laws.

In the simple one-dimensional case, suppose random variables $X_n$ converge weakly to $X$, so that their distribution functions satisfy
\[
  F_{X_n}(x)\to F_X(x)
\]
at all continuity points of $F_X$.

Skorokhod tells us that one can pass to a new probability space, for example the unit interval $(0,1)$ equipped with Lebesgue measure $\lambda$, and define new random variables by the quantile representation
\[
  \widetilde X_n(r)=F_{X_n}^{-1}(r), \qquad r\in(0,1).
\]
Then
\[
  \widetilde X_n(r)\to \widetilde X(r)
\]
almost everywhere, while the new variables retain exactly the same laws as the original ones. Indeed, for every $x\in\mathbb{R}$,
\[
  \lambda(\widetilde X_n\le x)
  =\lambda\bigl(\{r:F_{X_n}^{-1}(r)\le x\}\bigr)
  =\lambda\bigl(\{r:r\le F_{X_n}(x)\}\bigr)
  =F_{X_n}(x).
\]
So the price of gaining almost sure convergence is that we move to a new probability space.

\section*{7. The Paradigm Shift: Martingale Solutions}
This leads to the conceptual shift emphasized in the notes: after applying Skorokhod representation, we have transplanted the distributions of the approximate solutions onto a new probability space. In doing so, we lose the original filtration and the original driving Wiener process.

Therefore, we can no longer work with a strong probabilistic solution tied to a fixed Brownian motion. Instead, we pass to a weak probabilistic formulation, called a martingale solution.

A probability measure $P$ on $(\Omega,\mathcal F)$ is called a martingale solution to the SPDE
\[
  du(t)=A(u(t))\,dt+B(u(t))\,dW(t)
\]
if the following hold:
\begin{enumerate}
\item the initial condition is satisfied almost surely,
\[
  P\bigl(u(0)=u_0\bigr)=1;
\]
\item the process
\[
  M_t:=u(t)-u(0)-\int_0^t A(u(s))\,ds
\]
is a continuous $H$-valued martingale under $P$.
\end{enumerate}
Moreover, this martingale must have the quadratic variation prescribed by the diffusion coefficient:
\[
  \langle M\rangle_t=\int_0^t B(u(s))\,Q\,B(u(s))^*\,ds.
\]
Here $Q$ is the covariance operator on the noise space $U$, and $B(u(s))^*:H\to U$ denotes the adjoint operator.

This is what is meant here by a \emph{weak probabilistic form}. In a \emph{strong probabilistic form}, one fixes the probability space and the Wiener process first, and then solves the equation on that fixed stochastic basis. In a \emph{weak probabilistic form}, one is allowed to change the probability space and construct a new Wiener process together with the solution.

\subsection*{7.1. Detour: It\^o Representation Theorem}
To recover a stochastic integral formulation from the martingale characterized above, one uses the \textbf{It\^o representation theorem} in the form appropriate for continuous square-integrable martingales.

\begin{theorem}[It\^o representation theorem]
Let $M=(M_t)_{t\in[0,T]}$ be a martingale with continuous sample paths and finite second moment. Denote by $\langle M\rangle_t$ its quadratic variation, that is, the process measuring the accumulated variance of $M$ up to time $t$. Suppose that
\[
  \langle M\rangle_t=\int_0^t G(s)\,Q\,G(s)^*\,ds,
\]
where $Q$ is the covariance operator of the noise on a Hilbert space $U$, and $G(s)$ describes how the noise acts on the system at time $s$. Then there exists a Wiener process $\widetilde W$ such that
\[
  M_t=\int_0^t G(s)\,d\widetilde W(s),
  \qquad t\in[0,T].
\]
\end{theorem}

Because $M_t$ is a continuous martingale with this specific quadratic variation, the It\^o representation theorem allows one to reconstruct a new Wiener process $\widetilde W$ on the new probability space such that
\[
  M_t=\int_0^t B(u(s))\,d\widetilde W(s).
\]
So the SPDE is recovered in weak probabilistic form.

\section*{8. The Equivalent Formulation via Test Functions}
To actually prove that the limit of the Galerkin approximations satisfies the martingale property, one uses the equivalent formulation written at the very bottom of the page. This is the classical Stroock--Varadhan martingale problem approach, projected down to one dimension by testing against the basis vectors of $H$.

Fix a basis vector $e_k$ and a test function
\[
  \varphi\in C_b^2(\mathbb{R}).
\]
Instead of working directly with the abstract $H$-valued martingale $M_t$, we apply $\varphi$ to the scalar Fourier coefficient $\langle u(t),e_k\rangle_H$.

Define the scalar process
\[
  M_t^{k,\varphi}:=\varphi\bigl(\langle u(t),e_k\rangle_H\bigr)-\varphi\bigl(\langle u_0,e_k\rangle_H\bigr)
\]
\[
  \qquad -\int_0^t \varphi'\bigl(\langle u(s),e_k\rangle_H\bigr)\,\langle A(u(s)),e_k\rangle\,ds
\]
\[
  \qquad +\frac12\int_0^t \varphi''\bigl(\langle u(s),e_k\rangle_H\bigr)\,
  \bigl\langle (B(u(s))\,Q\,B(u(s))^*)e_k,e_k\bigr\rangle_H\,ds.
\]
This last term is the quadratic variation contribution of the diffusion, projected onto the basis vector $e_k$. The martingale problem formulation says that $M_t^{k,\varphi}$ is a martingale, equivalently,
\[
  \mathbb E_P\bigl[M_t^{k,\varphi}-M_s^{k,\varphi}\mid \mathcal F_s\bigr]=0
  \qquad \text{for } 0\le s\le t\le T.
\]
Equivalently, for every bounded continuous $\mathcal F_s$-measurable random variable $\Phi_s$,
\[
  \mathbb E_P\bigl[(M_t^{k,\varphi}-M_s^{k,\varphi})\Phi_s\bigr]=0.
\]
The same identity holds for the finite-dimensional laws $P_n$:
\[
  \mathbb E_{P_n}\bigl[(M_t^{k,\varphi}-M_s^{k,\varphi})\Phi_s\bigr]=0.
\]
Therefore, to prove that the limit law $P$ is a martingale solution, it is enough to verify these scalar identities for all $\varphi$, $e_k$, and $\Phi_s$. This is the point of the section: an abstract $H$-valued martingale statement is difficult to pass to the limit directly, while these scalar expectation identities are concrete and stable under weak convergence of the laws.

This page marks the functional shift of the lecture: the Minty--Browder monotonicity trick is replaced by compact embedding, Prokhorov compactness, and the Skorokhod representation theorem.

\section*{9. The Existence Proof via Tightness and Weak Convergence}
We now continue to the next page of the notes, where the martingale problem formulation is completed and the topological proof of existence is finished. This page combines It\^o's formula, Galerkin approximation, tightness, Prokhorov compactness, and weak convergence of probability laws into one single argument.
\subsection*{9.1. The Main Existence Theorem}
We can now state the central existence result in the martingale formulation.

\begin{theorem}[Theorem 3.13]
Under assumptions (V.1), (V.3), and (V.5)--(V.7), there exists a solution $P$ to the martingale problem associated with
\[
  du(t)=A(u(t))\,dt+B(u(t))\,dW(t).
\]
\end{theorem}

The decisive new ingredient is assumption (V.6), namely the compact embedding
\[
  V\hookrightarrow\hookrightarrow H.
\]
Because monotonicity has been dropped, one can no longer expect the previous variational argument to produce a unique strong solution on a fixed stochastic basis. What this theorem guarantees instead is the existence of a weak probabilistic solution, that is, a martingale solution.
In other words, the martingale problem is to find a probability law $P$ on the path space such that the canonical process satisfies the initial condition and the associated drift-corrected test processes are martingales under $P$.

\begin{proof}

\subsection*{9.2. Laws of the Galerkin Approximations}
To prove the theorem, we return to the Galerkin approximation, exactly as in the monotone case. We project the SPDE onto the finite-dimensional subspace
\[
  V_n=\mathrm{span}\{e_1,\dots,e_n\}.
\]
But now the focus shifts from the random variables themselves to their probability laws. Let
\[
  P_n=\mathrm{Law}(u_n)
\]
denote the law of the solution of the approximate finite-dimensional SDE.

For each $n$, the measure $P_n$ satisfies the following three properties:
\begin{enumerate}
\item its support is contained in the continuous path space over the Galerkin subspace,
\[
  \mathrm{supp}(P_n)\subset C([0,T];V_n);
\]
\item it satisfies the projected initial condition almost surely,
\[
  P_n\bigl(u(0)=\pi_n u_0\bigr)=1;
\]
\item it solves the finite-dimensional martingale problem, meaning that for every basis vector $e_k$, every $\varphi\in C_b^2(\mathbb{R})$, and every bounded past-dependent test random variable $\Phi_s$,
\[
  \mathbb E_{P_n}\bigl[(M_t^{k,\varphi}-M_s^{k,\varphi})\Phi_s\bigr]=0.
\]
\end{enumerate}

\subsection*{9.3. Tightness and the Role of Compact Embedding}
To pass to the limit as $n\to\infty$, we need compactness at the level of laws. This is where Prokhorov's theorem enters: once tightness has been established, it gives relative compactness of the family $(P_n)$ in the topology of weak convergence. In particular, one may extract a subsequence, still denoted by $(P_n)$, and a probability measure $P$ such that
\[
  P_n\Rightarrow P.
\]

\begin{lemma}[Lemma 3.14]
The sequence of probability measures $(P_n)_{n\ge 1}$ is tight on
\[
  \Omega:=C([0,T];V'_{\mathrm{weak}})\cap L^2([0,T];V)\cap L^2([0,T];H)
\]
with respect to the three topologies
\[
  \tau_1,\tau_2,\tau_3,
\]
where $\tau_1$ is the weak topology on $L^2([0,T];V)$, $\tau_2$ is the uniform topology on $C([0,T];V'_{\mathrm{weak}})$, and $\tau_3$ is the strong topology on $L^2([0,T];H)$.
\end{lemma}

The topology $\tau_1$ comes from the uniform a priori bounds in $L^2([0,T];V)$. The topology $\tau_2$ controls continuity of the trajectories in the dual space $V'$. The crucial point is $\tau_3$, the strong topology on $L^2([0,T];H)$, because strong compactness in $H$ is obtained precisely from the compact embedding
\[
  V\hookrightarrow\hookrightarrow H.
\]
Thus, control in the more regular space $V$ yields compactness in the stronger space $H$, and this is the key topological mechanism replacing monotonicity in the present lecture.

\subsection*{9.4. Passing to the Limit}
Once Lemma 3.14 is available, Prokhorov's theorem yields a subsequence, still denoted by $(P_n)$, and a probability measure $P$ such that
\[
  P_n\Rightarrow P.
\]
This measure $P$ is the candidate solution to the martingale problem.

To finish the proof, one must pass to the limit in the scalar martingale identity
\[
  \mathbb E_{P_n}\bigl[(M_t^{k,\varphi}-M_s^{k,\varphi})\Phi_s\bigr]=0.
\]
Weak convergence of the laws alone is not enough for this step. One must first establish uniform integrability of the martingale increments, and this is the content of the final page of the argument.

\section*{10. Final Step: Uniform Integrability of the Martingale Increments}
Welcome back. Let us seamlessly continue our lecture by examining this final page of notes. We are now at the last hurdle in the existence proof for the martingale solution.

Recall where we left off: via Prokhorov's theorem and the Skorokhod representation argument, we extracted a sequence of probability measures $P_n$ that converges weakly to a limit measure $P$. However, to prove that this limit measure $P$ is indeed a martingale solution, we must pass to the limit $n\to\infty$ inside the expectation of the martingale increments.

Weak convergence alone is not enough for this passage to the limit. One also needs uniform integrability. To obtain uniform integrability, it is enough to prove that the $p$-th moments of the martingale increments are uniformly bounded for some $p>1$.

\subsection*{10.1. The Goal: Bounding the Martingale Increments}
At the top of the page, one writes again the explicit definition of the scalar martingale process $M_t^{k,\varphi}$ associated with the sequence $u_n(t)$. One also recalls the continuity assumption on the drift,
\[
  A:V_{\mathrm{weak}}\cap H\to V'_{\mathrm{weak}}.
\]
The strategy is to combine these assumptions on the behavior of the operators, namely (V.1), (V.3), and (V.5)--(V.7), with the a priori estimate
\[
  \sup_n \mathbb E_{P_n}\Bigl(\sup_{t\le T}\|u_n(t)\|_H^2+\int_0^T \|u_n(s)\|_V^2\,ds\Bigr)<\infty
\]
in order to prove that for some $p>1$,
\[
  \sup_n \mathbb E_{P_n}\bigl(|M_t^{k,\varphi}-M_s^{k,\varphi}|^p\bigr)<\infty.
\]

Recall that
\[
  M_t^{k,\varphi}
  :=\varphi\bigl(\langle u(t),e_k\rangle_H\bigr)-\varphi\bigl(\langle u_0,e_k\rangle_H\bigr)
\]
\[
  \qquad -\int_0^t \varphi'\bigl(\langle u(r),e_k\rangle_H\bigr)\,\langle A(u(r)),e_k\rangle\,dr
\]
\[
  \qquad +\frac12\int_0^t \varphi''\bigl(\langle u(r),e_k\rangle_H\bigr)\,
  \bigl\langle (B(u(r))Q B(u(r))^*)e_k,e_k\bigr\rangle_H\,dr.
\]

\subsection*{10.2. Decomposing the Increment}
We break the increment $M_t^{k,\varphi}-M_s^{k,\varphi}$ into its three constituent parts.

The first part is the test-function term
\[
  \varphi\bigl(\langle u(t),e_k\rangle_H\bigr)-\varphi\bigl(\langle u(s),e_k\rangle_H\bigr).
\]
Since $\varphi\in C_b^2(\mathbb R)$, this term is bounded.

The second part is the drift integral
\[
  \int_s^t \varphi'\bigl(\langle u(r),e_k\rangle_H\bigr)\,\langle A(u(r)),e_k\rangle\,dr.
\]
Because $\varphi'$ is bounded and assumption (V.3) gives
\[
  \|A(u)\|_{V'}\le M(1+\|u\|_V),
\]
this term is controlled by a constant multiple of
\[
  \int_s^t (1+\|u(r)\|_V)\,dr.
\]

The third part is the second-order It\^o correction term
\[
  \frac12\int_s^t \varphi''\bigl(\langle u(r),e_k\rangle_H\bigr)\,
  \bigl\langle (B(u(r))Q B(u(r))^*)e_k,e_k\bigr\rangle_H\,dr.
\]
Again $\varphi''$ is bounded. To control the diffusion, one uses assumption (V.5),
\[
  \|B(u)\circ \sqrt Q\|_{L_2(U,H)}\le M(1+\|u\|_V^{1-\delta})
\]
for some $\delta>0$. Since the It\^o correction involves the square of the diffusion, this term is bounded by a constant multiple of
\[
  \int_s^t \bigl(1+\|u(r)\|_V^{2(1-\delta)}\bigr)\,dr.
\]

\subsection*{10.3. The Role of Sub-linear Growth}
Putting the three bounds together, one finds that the absolute increment is controlled by an expression of the form
\[
  |M_t^{k,\varphi}-M_s^{k,\varphi}|
  \le C\Bigl(1+\int_s^t \|u(r)\|_V^2\,dr\Bigr)^{1-\delta}.
\]
Here one uses H\"older's inequality and the fact that $1-\delta<1$:
\[
  \int_s^t \|u(r)\|_V^{2(1-\delta)}\,dr \le C\Bigl(\int_s^t \|u(r)\|_V^2\,dr\Bigr)^{1-\delta}.
\]
The critical feature is that the exponent $1-\delta$ lies outside the integral.

To obtain a uniform $L^p$-bound, choose
\[
  p=\frac{1}{1-\delta}>1.
\]
Then
\[
  p(1-\delta)=1.
\]
Therefore, taking the $p$-th power gives
\[
  |M_t^{k,\varphi}-M_s^{k,\varphi}|^p
  \le
  C\Bigl(1+\int_s^t \|u(r)\|_V^2\,dr\Bigr).
\]
Taking expectation under $P_n$ yields
\[
  \mathbb E_{P_n}\bigl(|M_t^{k,\varphi}-M_s^{k,\varphi}|^p\bigr)
  \le
  C\,\mathbb E_{P_n}\Bigl(1+\int_s^t \|u(r)\|_V^2\,dr\Bigr).
\]
Since
\[
  \int_s^t \|u(r)\|_V^2\,dr\le \int_0^T \|u(r)\|_V^2\,dr,
\]
the a priori estimate implies
\[
  \sup_n \mathbb E_{P_n}\bigl(|M_t^{k,\varphi}-M_s^{k,\varphi}|^p\bigr)<\infty.
\]
Hence the martingale increments are uniformly integrable.

\subsection*{10.4. The Final Conclusion}
This uniform integrability is the last ingredient needed to pass to the limit inside the expectation. Combined with the weak convergence obtained from Prokhorov's theorem, it yields
\[
  \lim_{n\to\infty}\mathbb E_{P_n}\bigl[(M_t^{k,\varphi}-M_s^{k,\varphi})\Phi_s\bigr]
  =
  \mathbb E_{P}\bigl[(M_t^{k,\varphi}-M_s^{k,\varphi})\Phi_s\bigr]
  =0.
\]
Therefore, under the limit measure $P$, the martingale increments still have zero conditional expectation given the past. This proves that $P$ is indeed a martingale solution to the infinite-dimensional SPDE.

The proof of Theorem 3.13 is complete.

\end{proof}

\section*{Appendix: Detailed Estimates for the Drift and Diffusion Terms}
For completeness, we record the detailed estimates used in the proof when bounding the scalar martingale increment.

\subsection*{A.1. Estimate for the Drift Term}
Consider the drift term
\[
  \int_s^t \varphi'\bigl(\langle u(r),e_k\rangle_H\bigr)\,\langle A(u(r)),e_k\rangle\,dr.
\]
We estimate the pairing $\langle A(u(r)),e_k\rangle$ by the dual norm of $A(u(r))$.

Since $\varphi\in C_b^2(\mathbb R)$, its derivative is bounded, so
\[
  \bigl|\varphi'\bigl(\langle u(r),e_k\rangle_H\bigr)\bigr|\le C_\varphi.
\]
Moreover, because $A(u(r))\in V'$ and $e_k\in V$, the duality estimate gives
\[
  |\langle A(u(r)),e_k\rangle|
  \le \|A(u(r))\|_{V'}\,\|e_k\|_V.
\]
Assumption (V.3) implies
\[
  \|A(u(r))\|_{V'}\le M(1+\|u(r)\|_V).
\]
Combining these bounds, we obtain
\[
  \bigl|\varphi'\bigl(\langle u(r),e_k\rangle_H\bigr)\,\langle A(u(r)),e_k\rangle\bigr|
  \le C_\varphi\,M\,\|e_k\|_V\,(1+\|u(r)\|_V).
\]
Since $C_\varphi$, $M$, and $\|e_k\|_V$ are constants, integration over $r\in[s,t]$ yields
\[
  \left|
  \int_s^t \varphi'\bigl(\langle u(r),e_k\rangle_H\bigr)\,\langle A(u(r)),e_k\rangle\,dr
  \right|
  \le C\int_s^t (1+\|u(r)\|_V)\,dr.
\]
This is why the drift contribution is controlled by a constant multiple of
\[
  \int_s^t (1+\|u(r)\|_V)\,dr.
\]

\subsection*{A.2. Estimate for the Diffusion Term}
Now consider the second-order It\^o correction term
\[
  \frac12\int_s^t \varphi''\bigl(\langle u(r),e_k\rangle_H\bigr)\,
  \bigl\langle (B(u(r))Q B(u(r))^*)e_k,e_k\bigr\rangle_H\,dr.
\]
Again $\varphi''$ is bounded, since $\varphi\in C_b^2(\mathbb R)$. It therefore remains to estimate
\[
  \bigl|\bigl\langle (B(u)Q B(u)^*)e_k,e_k\bigr\rangle_H\bigr|.
\]

Write this using $\sqrt Q$ by setting
\[
  T=B(u)\circ \sqrt Q.
\]
Then
\[
  TT^*=(B(u)\sqrt Q)(B(u)\sqrt Q)^*=B(u)Q B(u)^*.
\]
Hence
\[
  \bigl\langle (B(u)Q B(u)^*)e_k,e_k\bigr\rangle_H
  =
  \langle TT^*e_k,e_k\rangle_H
  =
  \|T^*e_k\|_U^2
  =
  \|(B(u)\sqrt Q)^*e_k\|_U^2.
\]
Using the operator norm estimate by the Hilbert--Schmidt norm, we get
\[
  \|(B(u)\sqrt Q)^*e_k\|_U
  \le \|B(u)\sqrt Q\|_{L_2(U,H)}\,\|e_k\|_H.
\]
Therefore
\[
  \bigl\langle (B(u)Q B(u)^*)e_k,e_k\bigr\rangle_H
  \le \|B(u)\sqrt Q\|_{L_2(U,H)}^2\,\|e_k\|_H^2.
\]
Now apply assumption (V.5):
\[
  \|B(u)\circ \sqrt Q\|_{L_2(U,H)}
  \le M(1+\|u\|_V^{1-\delta}).
\]
Squaring this inequality yields, after adjusting the constant,
\[
  \|B(u)\sqrt Q\|_{L_2(U,H)}^2
  \le C\bigl(1+\|u\|_V^{2(1-\delta)}\bigr).
\]
Substituting back, we arrive at
\[
  \bigl|\varphi''\bigl(\langle u(r),e_k\rangle_H\bigr)\,
  \langle (B(u(r))Q B(u(r))^*)e_k,e_k\rangle_H\bigr|
  \le C\bigl(1+\|u(r)\|_V^{2(1-\delta)}\bigr).
\]
After integration over $[s,t]$, this gives
\[
  \left|
  \frac12\int_s^t \varphi''\bigl(\langle u(r),e_k\rangle_H\bigr)\,
  \langle (B(u(r))Q B(u(r))^*)e_k,e_k\rangle_H\,dr
  \right|
  \le C\int_s^t \bigl(1+\|u(r)\|_V^{2(1-\delta)}\bigr)\,dr.
\]
This is why the diffusion contribution is controlled by a constant multiple of
\[
  \int_s^t \bigl(1+\|u(r)\|_V^{2(1-\delta)}\bigr)\,dr.
\]

\section*{11. Final Passage to the Limit by Truncation}
Excellent. Let us close the lecture by making explicit the final truncation argument that converts the uniform $L^p$-bound into convergence of the martingale expectations under the weakly convergent laws.

At this stage we already know that $P_n\Rightarrow P$. The remaining task is to prove the second martingale-problem condition under the limit measure $P$, namely that the martingale increments still have zero conditional expectation.

\subsection*{11.1. The Truncation Function}
Weak convergence alone allows one to pass to the limit only for bounded continuous test functions. The increment
\[
  M_t^{k,\varphi}-M_s^{k,\varphi}
\]
is continuous, but not automatically bounded. To force boundedness, fix $K>0$ and define the truncation
\[
  \Psi_K:=(-K)\vee (M_t^{k,\varphi}-M_s^{k,\varphi})\wedge K.
\]
Equivalently,
\[
  \Psi_K=
  \begin{cases}
    K, & M_t^{k,\varphi}-M_s^{k,\varphi}>K,\\
    M_t^{k,\varphi}-M_s^{k,\varphi}, & |M_t^{k,\varphi}-M_s^{k,\varphi}|\le K,\\
    -K, & M_t^{k,\varphi}-M_s^{k,\varphi}<-K.
  \end{cases}
\]
Thus $\Psi_K$ agrees with the martingale increment on $[-K,K]$ and is cut off outside this interval. In particular, $\Psi_K$ is bounded and continuous, and $\Psi_K\to M_t^{k,\varphi}-M_s^{k,\varphi}$ pointwise as $K\to\infty$.

Since $\Psi_K$ is bounded and continuous, and $\Phi_s$ is also bounded and continuous, weak convergence gives
\[
  \mathbb E_{P_n}(\Psi_K\Phi_s)\to \mathbb E_P(\Psi_K\Phi_s)
  \qquad \text{for every fixed } K.
\]

\subsection*{11.2. The Truncation Error}
We now estimate the error introduced by replacing the true increment by its truncated version. One has
\[
  \left|
  \mathbb E_{P_n}\Bigl[\bigl(\Psi_K-(M_t^{k,\varphi}-M_s^{k,\varphi})\bigr)\Phi_s\Bigr]
  \right|
  \le
  \|\Phi_s\|_\infty\,
  \mathbb E_{P_n}\Bigl(
  \bigl|\Psi_K-(M_t^{k,\varphi}-M_s^{k,\varphi})\bigr|
  \Bigr).
\]
This follows from the elementary inequalities $|\mathbb E[Y]|\le \mathbb E[|Y|]$ and $|Y\Phi_s|\le \|\Phi_s\|_\infty |Y|$.
The difference is nonzero only on the tail event
\[
  \{|M_t^{k,\varphi}-M_s^{k,\varphi}|>K\}.
\]
Indeed, if
\[
  |M_t^{k,\varphi}-M_s^{k,\varphi}|\le K,
\]
then the truncation does nothing and the difference is zero. On the tail event
\[
  |M_t^{k,\varphi}-M_s^{k,\varphi}|>K,
\]
one has
\[
  \bigl|\Psi_K-(M_t^{k,\varphi}-M_s^{k,\varphi})\bigr|
  \le
  |M_t^{k,\varphi}-M_s^{k,\varphi}|.
\]
Therefore
\[
  \left|
  \mathbb E_{P_n}\Bigl[\bigl(\Psi_K-(M_t^{k,\varphi}-M_s^{k,\varphi})\bigr)\Phi_s\Bigr]
  \right|
  \le
  \|\Phi_s\|_\infty\,
  \mathbb E_{P_n}\Bigl(
  |M_t^{k,\varphi}-M_s^{k,\varphi}|
  \mathbf 1_{\{|M_t^{k,\varphi}-M_s^{k,\varphi}|>K\}}
  \Bigr).
\]

\subsection*{11.3. Using the Uniform $L^p$-Bound}
Now use the $p$-th moment estimate proved in the previous section. On the event
\[
  |M_t^{k,\varphi}-M_s^{k,\varphi}|>K,
\]
we have
\[
  |M_t^{k,\varphi}-M_s^{k,\varphi}|
  \le
  \frac{|M_t^{k,\varphi}-M_s^{k,\varphi}|^p}{K^{p-1}}.
\]
Hence
\[
  \mathbb E_{P_n}\Bigl(
  |M_t^{k,\varphi}-M_s^{k,\varphi}|
  \mathbf 1_{\{|M_t^{k,\varphi}-M_s^{k,\varphi}|>K\}}
  \Bigr)
  \le
  \frac{1}{K^{p-1}}
  \mathbb E_{P_n}\bigl(|M_t^{k,\varphi}-M_s^{k,\varphi}|^p\bigr).
\]
Since $p>1$, the factor $K^{-(p-1)}$ tends to zero as $K\to\infty$. Because the $p$-th moments are uniformly bounded in $n$, the truncation error converges to zero uniformly in $n$, and this uniform control is exactly what allows one to pass from the truncated variables to the original martingale increment in the limit.

\subsection*{11.4. Conclusion}
As a consequence of Section 11.3, one may first pass to the limit for the truncated variables and then let $K\to\infty$ to recover the original increment. Hence
\[
  \lim_{n\to\infty}
  \mathbb E_{P_n}\bigl[(M_t^{k,\varphi}-M_s^{k,\varphi})\Phi_s\bigr]
  =
  \mathbb E_P\bigl[(M_t^{k,\varphi}-M_s^{k,\varphi})\Phi_s\bigr].
\]
Since the left-hand side is zero for every $n$, it follows that
\[
  \mathbb E_P\bigl[(M_t^{k,\varphi}-M_s^{k,\varphi})\Phi_s\bigr]=0.
\]
This is exactly the desired martingale identity under the limit measure $P$.

\end{document}
